% ---------------------------------------------------------------------------
% Standalone TikZ source, recolored/adapted from the paper's own Figure 1
% (content/paper-PNAS-final/fig01_shannon_model_extension_standalone.tex) --
% the "extension" row. Unlike the paper's own figure (which continues
% rightward: logic sentence -> f -> g -> ...deductive machinery -> deduced
% facts, all in one wide horizontal line), this version keeps only the
% four stages shared with the top row horizontal -- using the *exact same*
% matrix column-sep/node styles as shannon_top_row.tex, so that shared
% prefix is pixel-identical to the top row's own message->f->g->recovered
% message once both are rendered at the same scale -- then turns the two
% genuinely new stages (deductive machinery, deduced facts) *downward*
% instead of further rightward. This keeps the whole diagram's own width
% close to the top row's own (rather than roughly double it), which is
% what actually didn't fit on a phone screen in the first version.
% Compiles to shannon_bottom_row.pdf, then converted to
% shannon_bottom_row.svg via pdftocairo -svg (see README.md in this
% directory for the exact build commands).
% ---------------------------------------------------------------------------
\documentclass[border=2pt]{standalone}

\usepackage[utf8]{inputenc}
\usepackage[T1]{fontenc}
% Sans-serif, matching the legend card's own font stack (-apple-system,
% ..., Helvetica, Arial, sans-serif) -- LaTeX's own default (Computer
% Modern) is a serif face, visibly mismatched against the HTML card
% sitting right next to these diagrams. `helvet` is a standard, always-
% bundled psnfss package (no extra font install needed).
\usepackage{helvet}
\renewcommand{\familydefault}{\sfdefault}
\usepackage{tikz}

\usetikzlibrary{shapes.geometric}
\usetikzlibrary{shapes.symbols}
\usetikzlibrary{dsp}
\usetikzlibrary{chains}
\usetikzlibrary{arrows,chains,matrix,positioning,scopes,patterns,decorations.pathreplacing,calc,positioning,fit,decorations.pathreplacing,calligraphy}

\definecolor{classical}{RGB}{184,255,250}
\definecolor{newstage}{RGB}{232,162,59}

\begin{document}
\normalfont
\begin{tikzpicture}[color=classical, text=classical]

	% Identical column-sep/row-sep and node styles to shannon_top_row.tex's
	% own single row -- this is what guarantees the shared prefix lands at
	% the exact same physical (pt) positions as the top row's own, not
	% merely a visually-similar approximation.
	\matrix[row sep=2.5mm, column sep=5mm](A)
	{	tm

		\node[dspnodeopen,dsp/label=left] (mm00) {logic sentence};  &
		\node[dspsquare]         (mm01) {f};     &
            \node[coordinate,label=bits] (mm02) {}; &
		\node[dspsquare]          (mm03) {g};     &
		\node[dspnodeopen,dsp/label=right]          (mm03p5) {recovered logic sentence};     &
  \\
  };
  \node[anchor=east, text opacity=0, xshift=-\dspoperatorlabelspacing] at (mm00.west) {logic sentence};
  \node[anchor=west, text opacity=0, xshift=\dspoperatorlabelspacing] at (mm03p5.east) {recovered logic sentence};

	\begin{scope}[start chain]
		\chainin (mm00);
		\chainin (mm01) [join=by dspconn];
		\chainin (mm03) [join=by dspconn];
		\chainin (mm03p5) [join=by dspconn];
        \end{scope}

  % ---- The genuinely new part: turns downward from here, rather than
  % continuing rightward -- so the whole diagram's own width stays close
  % to the top row's own, using vertical (scroll-friendly) space instead.
  \node[dspcloud, below=9mm of mm03p5, color=newstage, text=newstage] (machinery) {deductive\\machinery};
  \draw[dspconn, draw=newstage] (mm03p5.south) -- (machinery.north);

  \node[dspnodeopen, right=14mm of machinery, color=newstage, text=newstage, dsp/label=below] (facts) {deduced facts};
  \draw[dspconn, draw=newstage] (machinery.east) -- (facts.west);
  \node[anchor=north, text opacity=0, yshift=-\dspoperatorlabelspacing] at (facts.south) {deduced facts};

\end{tikzpicture}
\end{document}
